\documentclass[11pt]{article}

\usepackage[utf8]{inputenc}
\usepackage[T1]{fontenc}
\usepackage{lmodern}
\usepackage{geometry}
\usepackage{amsmath,amssymb,amsfonts,amsthm,mathtools}
\usepackage{bm}
\usepackage{enumitem}
\usepackage{microtype}
\usepackage{hyperref}
\usepackage{titlesec}
\usepackage{booktabs}
\usepackage{array}
\usepackage{setspace}
\usepackage{mathrsfs}
\usepackage{physics}

\geometry{margin=1in}
\hypersetup{colorlinks=true, linkcolor=black, urlcolor=blue, citecolor=black}
\setlist[enumerate]{topsep=0.4em,itemsep=0.35em}
\setlist[itemize]{topsep=0.3em,itemsep=0.25em}
\titleformat{\section}{\large\bfseries}{\thesection.}{0.5em}{}
\titleformat{\subsection}{\normalsize\bfseries}{\thesubsection.}{0.5em}{}
\setstretch{1.06}

\newcommand{\Aether}{\AE{}ther}
\newcommand{\Aflow}{\AE{}ther-flow}
\newcommand{\TheoryName}{\Aflow{} Interpretation of Relativity}
\newcommand{\TheoryProgram}{\Aether{} / \Aflow{} framework}
\newcommand{\Sub}{\mathcal A}
\newcommand{\ProjSub}{P}
\newcommand{\SpatialD}{\mathcal D}
\newcommand{\LapPerp}{\Delta_\perp}
\newcommand{\ObserverMetric}{g^{\mathrm{br}}}
\newcommand{\CandidateMetric}{g^{\mathrm{cand}}}
\newcommand{\BaseShift}{\Sigma^{\mathrm{IER}}}
\newcommand{\ResponseShift}{\Sigma^{\mathrm{resp}}}
\newcommand{\AuxPol}{\mathbb K}
\newcommand{\FlowPot}{\Phi_\Omega}
\newcommand{\SourceDensity}{\varrho_\Omega}
\newcommand{\SourceDensityRed}{\varrho_\Omega^{\mathrm{red}}}
\newcommand{\SourceMult}{\Lambda_\rho}
\newcommand{\EinLin}{\mathcal E_{\eta}}
\newcommand{\BaseAux}{\mathfrak X_{\mathrm{base}}}
\newcommand{\GaugeAux}{\mathfrak G_{\mathrm{gc}}}
\newcommand{\BaseEll}{\mathbb A_{\mathrm{base}}}
\newcommand{\FullLin}{\mathbb L_{\mathrm{full}}}
\newcommand{\RespLin}{\mathbb L_{\mathrm{resp}}}
\newcommand{\HamChargeOmega}{\mathfrak m_\Omega}
\newcommand{\MatterAct}{S_{\mathrm{matter}}}
\newcommand{\RespAct}{S_{\mathrm{resp}}}
\newcommand{\TotalAct}{S_{\mathrm{tot}}}
\newcommand{\RedAct}{S_{\mathrm{red}}}
\newcommand{\TT}{\mathrm{TT}}
\newcommand{\Ncand}{N_{\mathrm{cand}}}
\newcommand{\Gcand}{\gamma^{\mathrm{cand}}}

\newtheorem{definition}{Definition}
\newtheorem{proposition}{Proposition}
\newtheorem{remark}{Remark}
\newtheorem{corollary}{Corollary}
\newtheorem{theorem}{Theorem}

\title{\textbf{The \TheoryName{}}\\[0.4em]
\large Substrate-Response / Self-Response Charge-Polarization Candidate\\
\large Full Infrared-Spectrum and Universal-Coupling Gate Note}
\author{}
\date{}

\begin{document}

\maketitle

\begin{abstract}
The benchmark-facing charge-polarization candidate note already completed the first required step for the surviving substrate-response / self-response family: it wrote the sector, derived the eliminated observer equation on matter-free reduced data first, preserved one operative metric, and gave a first infrared screen. The present manuscript is the next benchmark-gate follow-up fixed by the live research plan. Its task is no longer family selection and no longer another deeper origin continuation. It is the Phase D and Phase E test for that concrete candidate.

The result is positive at the present disciplined scope. About an admissible homogeneous benchmark branch, the \emph{full} coupled reduced linearization factorizes into the already-admissible base GR-compatible block and an explicit charge-polarization response block. Because the reduced Hamiltonian source density satisfies \(\SourceDensityRed=\mathcal O_{\ge 2}\), the new response sector has no linear mixing with the base observer or auxiliary variables. Its linear symbol is purely elliptic or algebraic:
\[
\RespLin(\mathbf{k})
=
\begin{pmatrix}
0 & 1 & 0 \\
1 & 0 & -1 \\
0 & -1 & c_\Phi |\mathbf{k}|^2
\end{pmatrix}
\oplus \mu_t^2 \oplus \mu_s^2 \oplus \mu_v^2\mathbf 1_3 \oplus \mu_{\mathrm{tf}}^2\mathbf 1_5,
\qquad
c_\Phi>0.
\]
Accordingly, the observer-accessible infrared poles of the full coupled system are exactly the ordinary two tensor poles \(\omega^2=|\mathbf{k}|^2\) of the base branch. No unsuppressed scalar, vector, preferred-frame, ghost, tachyonic, or second-cone remainder is generated by the charge-polarization completion.

The note then re-derives the matter and light sector after elimination rather than relying only on architectural wording. Eliminating \((\SourceDensity,\SourceMult,\FlowPot,\AuxPol)\) changes the gravitational action by a charge-polarization functional but does not introduce any new matter argument. The reduced action remains
\[
\RedAct[\CandidateMetric,\psi]
=
S_{\mathrm{grav}}^{\mathrm{cand}}[\CandidateMetric;\text{reduced substrate data}]
+
\MatterAct[\CandidateMetric,\psi].
\]
Therefore matter and light continue to couple only through the single operative metric \(\CandidateMetric\). In particular, Maxwell geometric optics follows the \(\CandidateMetric\)-null cone, massive test bodies follow \(\CandidateMetric\)-geodesics, and redshift and proper time are recovered from the same \(\CandidateMetric\) lapse and line element.

This note still does not claim the controlled GR-limit theorem or the final residual-sector verdict. Its positive result is narrower and exactly the one now required: the concrete charge-polarization candidate passes the full infrared-spectrum gate and the universal matter/light-coupling gate at current admissible-branch scope, so only the controlled GR-limit and residual-sector classifications remain as the next primary derivational manuscripts.
\end{abstract}

\section{Introduction}

The exact-closure benchmark remains the public front door of the repository, and the benchmark-facing derivation-gates note remains the criterion by which every derivational continuation is judged. The observer-remainder classification note and the charge-polarization candidate note have already fixed the next bridge-level task more narrowly than before. The question is no longer whether one can write another candidate-family survey, another packaging note, or another anchored positive-pair continuation. The question is whether the recorded charge-polarization candidate survives the real Phase D and Phase E gates.

That question has two parts. First, one must linearize the \emph{full} coupled reduced system about an admissible homogeneous branch, not just the newly added response block, and derive the observer-accessible spectrum explicitly. Second, one must re-derive the matter and light sector \emph{after elimination} and verify that the single operative metric survives at the same scope for null cones, geodesics, redshift, and proper time.

\paragraph{Framework and claim boundary.}
\TheoryProgram{} denotes the broader \Aether{} / \Aflow{} framework built on the claims that the \Aether{} is the underlying four-dimensional substrate of reality, the \Aflow{} is its intrinsic ordered motion, observed three-dimensional space is the local experiential slice of that deeper substrate, S-time is the experienced order of change arising from matter, light, and the \Aflow{}, and observed expansion is the three-dimensional appearance of deeper four-dimensional ordered motion. Within that framework, \TheoryName{} denotes the exact-closure theory in which the effective gravitational dynamics are adopted to be exactly Einsteinian with universal matter coupling. In the active sequence, \emph{adoption} means use of that established relativistic dynamics without claiming substrate derivation, while \emph{derivation} is reserved for a first-principles recovery from explicit substrate variables.

The present note stays strictly inside that boundary. It does not claim that GR has already been derived from substrate microphysics. It does not promote the charge-polarization candidate to the controlled GR-limit theorem. It performs the narrower next benchmark-gate test now required: the full coupled infrared-spectrum verdict and the universal matter/light-coupling check for the concrete recorded candidate.

\section{Fixed Inputs and Benchmark-Branch Scope}

\begin{definition}[Fixed candidate inputs]
\label{def:inputs}
For the present note, the following are treated as fixed inputs.
\begin{enumerate}
    \item The benchmark exact-closure package and the derivation-gates note remain the fixed observer-level target.
    \item The charge-polarization candidate note fixes the response action, the eliminated response equations, the operative metric
    \begin{equation}
    \CandidateMetric_{\mu\nu}
    =
    \ObserverMetric_{\mu\nu}
    +\BaseShift_{\mu\nu}
    +\ResponseShift_{\mu\nu},
    \label{eq:gcand}
    \end{equation}
    and the fact that
    \begin{equation}
    \SourceDensityRed[Q,W,\chi]=\mathcal O_{\ge 2}
    \label{eq:rhoorder}
    \end{equation}
    on the active reduced-data counting.
    \item The same note already proves that the eliminated response shift is quadratic in \(\FlowPot\) and therefore begins only at quartic reduced-data order:
    \begin{equation}
    \AuxPol_{AB}^{\mathrm{elim}}
    =
    -2\FlowPot^2 U_AU_B
    +\frac32\FlowPot^2\ProjSub_{AB},
    \qquad
    \ResponseShift_{\mu\nu}=\mathcal O_{\ge 4}.
    \label{eq:resporder}
    \end{equation}
    \item The admissible base branch already used there is assumed to be one of the previously established GR-compatible linearized branches of the active higher-derivative line.
\end{enumerate}
\end{definition}

\begin{definition}[Admissible homogeneous benchmark branch]
\label{def:branch}
An \emph{admissible homogeneous benchmark branch} for the present note is a homogeneous background such that:
\begin{enumerate}
    \item
    \begin{equation}
    \overline{\CandidateMetric}_{\mu\nu}=\eta_{\mu\nu},
    \qquad
    \overline Q=\overline W=\overline\chi=0,
    \qquad
    \overline{\SourceDensity}
    =
    \overline{\SourceMult}
    =
    \overline{\FlowPot}
    =
    \overline{\AuxPol}_{AB}
    =0;
    \label{eq:branchdata}
    \end{equation}
    \item the ordered-flow background is homogeneous with fixed \(U^A\) and \(\ProjSub_{AB}\);
    \item after gauge and constraint extraction, the base reduced linearized operator takes the already-admissible GR-compatible form
    \begin{equation}
    \mathbb L_{\mathrm{base}}(\omega,\mathbf{k})
    =
    (\omega^2-|\mathbf{k}|^2)\mathbf 1_{\TT}
    \oplus
    \BaseEll(\mathbf{k}),
    \label{eq:baseoperator}
    \end{equation}
    where \(\mathbf 1_{\TT}\) acts on the two transverse-traceless graviton polarizations and \(\BaseEll(\mathbf{k})\) has no observer-accessible infrared pole:
    \begin{equation}
    \det \BaseEll(\mathbf{0})\neq 0.
    \label{eq:baseinvertible}
    \end{equation}
\end{enumerate}
\end{definition}

\begin{remark}
Definition \ref{def:branch} is the exact place where the present note uses the already-recorded base-branch admissibility rather than reopening it. The new task is to test the recorded charge-polarization completion on top of that branch, not to reprove the earlier base linearized-consistency route.
\end{remark}

\section{Full Coupled Linearization and the Infrared Spectrum}

\begin{definition}[Full linearized variable split]
\label{def:split}
Write the full linearized reduced state as
\begin{equation}
\delta\mathfrak U
=
\bigl(h^{\TT}_{ij},\GaugeAux,\BaseAux;
\delta\SourceMult,\delta\SourceDensity,\delta\FlowPot,
\delta\alpha,\delta\beta,\delta\mathcal V_i,\delta\mathcal T_{ij}\bigr),
\label{eq:split}
\end{equation}
where:
\begin{enumerate}
    \item \(h^{\TT}_{ij}\) denotes the observer-accessible transverse-traceless metric perturbation of \(\CandidateMetric\);
    \item \(\GaugeAux\) collects the gauge/constraint-exact metric variables;
    \item \(\BaseAux\) denotes the already-existing nonmetric base auxiliaries of the admissible branch;
    \item \((\delta\SourceMult,\delta\SourceDensity,\delta\FlowPot,\delta\alpha,\delta\beta,\delta\mathcal V_i,\delta\mathcal T_{ij})\) are the linearized charge-polarization variables, with \(\partial^i\delta\mathcal V_i=0\) and \(\delta^{ij}\delta\mathcal T_{ij}=0=\partial^i\delta\mathcal T_{ij}\) on the homogeneous slice.
\end{enumerate}
\end{definition}

\begin{proposition}[Quadratic charge-polarization block]
\label{prop:quadratic}
About an admissible homogeneous benchmark branch, the quadratic response action is
\begin{align}
\delta^2\RespAct
=
\int_{\Sub} d^4X\,\sqrt{\bar G}\,
\Bigl[
&\delta\SourceMult\,\delta\SourceDensity
-\delta\SourceDensity\,\delta\FlowPot
-\frac{1}{8\pi G_N}\,\partial_i\delta\FlowPot\,\partial_i\delta\FlowPot
\nonumber\\
&-\frac{\mu_t^2}{2}(\delta\alpha)^2
-\frac{\mu_s^2}{2}(\delta\beta)^2
-\frac{\mu_v^2}{2}\,\delta\mathcal V_i\delta\mathcal V_i
-\frac{\mu_{\mathrm{tf}}^2}{2}\,\delta\mathcal T_{ij}\delta\mathcal T_{ij}
\Bigr].
\label{eq:quadraticresp}
\end{align}
In particular, the response block has no linear mixing with \((h^{\TT}_{ij},\GaugeAux,\BaseAux)\).
\end{proposition}

\begin{proof}
Expand the response action of the candidate note around \eqref{eq:branchdata}. Because \(\SourceDensityRed=\mathcal O_{\ge 2}\) by \eqref{eq:rhoorder}, its linearization about the homogeneous branch vanishes. Because the eliminated polarization shift is quadratic in \(\FlowPot\), the terms \((\alpha+2\FlowPot^2)^2\) and \((\beta-\tfrac32\FlowPot^2)^2\) contribute only \((\delta\alpha)^2\) and \((\delta\beta)^2\) at quadratic order. No term linear in the response fields is multiplied by any base perturbation. This yields \eqref{eq:quadraticresp}.
\end{proof}

\begin{corollary}[Linearized response equations]
\label{cor:respeqs}
The linearized Euler--Lagrange equations of \eqref{eq:quadraticresp} are
\begin{equation}
\delta\SourceDensity=0,
\qquad
\delta\SourceMult=\delta\FlowPot,
\qquad
-\Delta\,\delta\FlowPot=0,
\qquad
\delta\alpha=\delta\beta=\delta\mathcal V_i=\delta\mathcal T_{ij}=0.
\label{eq:linearresp}
\end{equation}
On the decaying branch,
\begin{equation}
\delta\FlowPot=0
\qquad\Longrightarrow\qquad
\delta\ResponseShift_{\mu\nu}=0.
\label{eq:deltasigma}
\end{equation}
\end{corollary}

\begin{proof}
Variation of \eqref{eq:quadraticresp} with respect to \(\delta\SourceMult\) gives \(\delta\SourceDensity=0\). Variation with respect to \(\delta\SourceDensity\) then gives \(\delta\SourceMult=\delta\FlowPot\). The \(\delta\FlowPot\)-equation is the flat slice Laplace equation. Positive \(\mu_t^2,\mu_s^2,\mu_v^2,\mu_{\mathrm{tf}}^2\) imply the remaining algebraic equalities. On the decaying branch, the only harmonic \(\delta\FlowPot\) is zero, and \eqref{eq:resporder} then forces \(\delta\ResponseShift_{\mu\nu}=0\).
\end{proof}

\begin{proposition}[Full coupled linearized reduced system]
\label{prop:fullsystem}
After gauge and constraint extraction, the full coupled linearized reduced system about an admissible homogeneous benchmark branch is
\begin{align}
(\partial_0^2-\Delta)\,h^{\TT}_{ij} &= 0,
\label{eq:ttwave}
\\
\BaseEll(-\Delta)\,\BaseAux &= 0,
\label{eq:baseaux}
\\
\delta\SourceDensity &=0,
\qquad
\delta\SourceMult=\delta\FlowPot,
\qquad
-\Delta\,\delta\FlowPot=0,
\label{eq:phiblock}
\\
\delta\alpha&=\delta\beta=\delta\mathcal V_i=\delta\mathcal T_{ij}=0.
\label{eq:auxblock}
\end{align}
In particular, the full linearized operator is block diagonal:
\begin{equation}
\FullLin(\omega,\mathbf{k})
=
(\omega^2-|\mathbf{k}|^2)\mathbf 1_{\TT}
\oplus
\BaseEll(\mathbf{k})
\oplus
\RespLin(\mathbf{k}).
\label{eq:blockdiag}
\end{equation}
\end{proposition}

\begin{proof}
Definition \ref{def:branch} gives the base operator \eqref{eq:baseoperator}. Proposition \ref{prop:quadratic} shows that the response block has no linear mixing with the base variables, and Corollary \ref{cor:respeqs} gives its Euler--Lagrange equations explicitly. Combining the two yields \eqref{eq:ttwave}--\eqref{eq:blockdiag}.
\end{proof}

\begin{proposition}[Explicit response symbol and infrared poles]
\label{prop:symbol}
With the response variables ordered as
\[
(\delta\SourceMult,\delta\SourceDensity,\delta\FlowPot,
\delta\alpha,\delta\beta,\delta\mathcal V_i,\delta\mathcal T_{ij}),
\]
the response symbol takes the form
\begin{equation}
\RespLin(\mathbf{k})
=
\begin{pmatrix}
0 & 1 & 0 \\
1 & 0 & -1 \\
0 & -1 & c_\Phi |\mathbf{k}|^2
\end{pmatrix}
\oplus \mu_t^2 \oplus \mu_s^2 \oplus \mu_v^2\mathbf 1_3 \oplus \mu_{\mathrm{tf}}^2\mathbf 1_5,
\qquad
c_\Phi>0.
\label{eq:respsymbol}
\end{equation}
Hence
\begin{equation}
\det\RespLin(\mathbf{k})
=
-c_\Phi |\mathbf{k}|^2\,
\mu_t^2\mu_s^2(\mu_v^2)^3(\mu_{\mathrm{tf}}^2)^5,
\label{eq:respdet}
\end{equation}
so the response block introduces no \(\omega\)-dependent pole and no propagating infrared mode.
\end{proposition}

\begin{proof}
Equation \eqref{eq:quadraticresp} gives the response Hessian directly. The multiplier-density-potential block yields the \(3\times 3\) matrix in \eqref{eq:respsymbol} up to the positive Laplace coefficient \(c_\Phi\), while the polarization scalars, vector, and tracefree tensor contribute positive algebraic masses. Determinant expansion gives \eqref{eq:respdet}. Because \(\RespLin\) is independent of \(\omega\), it contributes no propagating observer-side pole.
\end{proof}

\begin{theorem}[Phase D infrared-spectrum verdict for the charge-polarization candidate]
\label{thm:phaseD}
On an admissible homogeneous benchmark branch, the full coupled reduced linearization of the charge-polarization candidate has the following observer-accessible spectrum.
\begin{enumerate}
    \item The only unsuppressed propagating observer-side poles are the ordinary two tensor poles
    \begin{equation}
    \omega^2=|\mathbf{k}|^2,
    \label{eq:tensorpole}
    \end{equation}
    carried by the two transverse-traceless graviton polarizations of \(h^{\TT}_{ij}\).
    \item Every base nonmetric sector remains in the already-admissible auxiliary or elliptic block \(\BaseEll(\mathbf{k})\), and every new charge-polarization variable belongs to the explicit elliptic or algebraic block \(\RespLin(\mathbf{k})\).
    \item No unsuppressed preferred-frame scalar, vector, ghost, tachyonic, or second-cone remainder survives at linear infrared order.
\end{enumerate}
\end{theorem}

\begin{proof}
By Proposition \ref{prop:fullsystem}, the full linearized operator factorizes as in \eqref{eq:blockdiag}. The base admissibility conditions \eqref{eq:baseoperator}--\eqref{eq:baseinvertible} say that the base branch already contributes only the Fierz--Pauli tensor pole together with an invertible nonpropagating block. Proposition \ref{prop:symbol} shows that the response block adds only elliptic or algebraic symbols and no new \(\omega\)-dependent pole.

Therefore the observer-accessible linear poles are exactly \eqref{eq:tensorpole}. Because the only propagating pole is the ordinary Fierz--Pauli tensor pole, the linear residue is the same positive tensor residue as on the admissible base branch, so no ghost is introduced. Because no additional pole of the form \(\omega^2=|\mathbf{k}|^2+m^2\) or \(\omega^2=|\mathbf{k}|^2-m^2\) appears, there is no unsuppressed tachyonic scalar or vector remainder. Because the response block is independent of \(\omega\) and because \(\delta\ResponseShift_{\mu\nu}=0\) by \eqref{eq:deltasigma}, the new sector does not create a second linear causal cone or an unsuppressed preferred-frame propagation law. This is exactly the required Phase D verdict.
\end{proof}

\begin{remark}
Theorem \ref{thm:phaseD} is stronger than the first infrared screen of the candidate note. That earlier note linearized only the added response block and then compared with the base branch. The present theorem linearizes the full coupled reduced system, writes the combined block structure explicitly, and identifies the observer-accessible poles of the whole candidate branch.
\end{remark}

\section{Matter and Light After Elimination}

\begin{definition}[Eliminated candidate action]
\label{def:redaction}
Let
\begin{equation}
\bigl(
\SourceDensity,\SourceMult,\FlowPot,\AuxPol
\bigr)
=
\bigl(
\SourceDensityRed,\FlowPot[\SourceDensityRed],\FlowPot[\SourceDensityRed],\AuxPol^{\mathrm{elim}}[\FlowPot]
\bigr)
\label{eq:elimsub}
\end{equation}
denote the eliminated charge-polarization solution of the candidate note. The \emph{eliminated candidate action} is the reduced action obtained by substituting \eqref{eq:elimsub} into the total action:
\begin{equation}
\RedAct
\equiv
\TotalAct\Big|_{(\SourceDensity,\SourceMult,\FlowPot,\AuxPol)=\mathrm{elim}}.
\label{eq:redaction}
\end{equation}
\end{definition}

\begin{proposition}[One metric survives elimination exactly]
\label{prop:one}
The eliminated candidate action has the form
\begin{equation}
\RedAct[\CandidateMetric,\psi]
=
S_{\mathrm{grav}}^{\mathrm{cand}}[\CandidateMetric;Q,W,\chi]
+
\MatterAct[\CandidateMetric,\psi],
\label{eq:redsplit}
\end{equation}
where \(S_{\mathrm{grav}}^{\mathrm{cand}}\) is independent of the matter fields \(\psi\). In particular, elimination introduces no second matter metric and no direct matter dependence on \(\SourceDensity\), \(\FlowPot\), or \(\AuxPol\).
\end{proposition}

\begin{proof}
The response action contains no matter field and no independent matter-coupling tensor. Its only observer-side imprint is the eliminated metric shift \(\ResponseShift\), hence the operative metric \(\CandidateMetric\). After substituting the eliminated solution \eqref{eq:elimsub}, every remaining response contribution is a functional of the reduced substrate variables and of \(\CandidateMetric\), but not of \(\psi\). Therefore the matter sector remains exactly \(\MatterAct[\CandidateMetric,\psi]\), establishing \eqref{eq:redsplit}.
\end{proof}

\begin{corollary}[No direct preferred-frame matter term after elimination]
\label{cor:nopf}
The eliminated candidate action contains no unsuppressed direct low-energy matter term of the forms
\begin{equation}
J_\mu U^\mu,
\qquad
\FlowPot\,\mathcal O_{\mathrm{matter}},
\qquad
\AuxPol_{AB}\,\mathcal O^{AB}_{\mathrm{matter}},
\label{eq:forbidden}
\end{equation}
outside the single metric dependence already encoded in \(\MatterAct[\CandidateMetric,\psi]\).
\end{corollary}

\begin{proof}
Proposition \ref{prop:one} shows that the matter fields enter only through \(\MatterAct[\CandidateMetric,\psi]\). Any term of the forms \eqref{eq:forbidden} would supply a new matter argument beyond \(\CandidateMetric\), which does not occur in \eqref{eq:redsplit}.
\end{proof}

\begin{proposition}[Matter equations after elimination]
\label{prop:mattereq}
For every matter species in the active exact-closure line, the eliminated matter equations are obtained by varying \(\MatterAct[\CandidateMetric,\psi]\) at fixed \(\CandidateMetric\):
\begin{equation}
\frac{\delta \MatterAct}{\delta\psi}
\biggr|_{\CandidateMetric}
=
0.
\label{eq:mattereom}
\end{equation}
Consequently, the principal symbol of the matter equations is determined solely by \(\CandidateMetric\).
\end{proposition}

\begin{proof}
In \eqref{eq:redsplit}, the only \(\psi\)-dependence of the eliminated action is through \(\MatterAct[\CandidateMetric,\psi]\). Hence variation with respect to \(\psi\) yields \eqref{eq:mattereom}. The principal symbol of the resulting matter equations is therefore the principal symbol of the corresponding minimally coupled matter system on \(\CandidateMetric\), with no extra contribution from eliminated response variables.
\end{proof}

\begin{theorem}[Null cone, geodesic, redshift, and proper-time recovery through \(\CandidateMetric\)]
\label{thm:phaseE}
At the eliminated-candidate scope, matter and light continue to recover the standard observer-side structures through the single operative metric \(\CandidateMetric\).
\begin{enumerate}
    \item \textbf{Null cone.} For Maxwell theory minimally coupled to \(\CandidateMetric\),
    \begin{equation}
    S_{\mathrm{Max}}[\CandidateMetric,A]
    =
    -\frac14
    \int d^4x\,\sqrt{-\CandidateMetric}\,
    g_{\mathrm{cand}}^{\mu\alpha}g_{\mathrm{cand}}^{\nu\beta}
    F_{\mu\nu}F_{\alpha\beta},
    \label{eq:maxwell}
    \end{equation}
    the geometric-optics eikonal satisfies
    \begin{equation}
    g_{\mathrm{cand}}^{\mu\nu}k_\mu k_\nu=0,
    \qquad
    k^\nu\nabla^{(\CandidateMetric)}_\nu k^\mu=0.
    \label{eq:nullcone}
    \end{equation}
    \item \textbf{Massive geodesics.} For a point particle,
    \begin{equation}
    S_{\mathrm{pp}}
    =
    -mc
    \int
    \sqrt{-\CandidateMetric_{\mu\nu}\dot x^\mu\dot x^\nu}\,d\lambda,
    \label{eq:pp}
    \end{equation}
    the Euler--Lagrange equation is the \(\CandidateMetric\)-geodesic equation
    \begin{equation}
    \ddot x^\mu+\Gamma^\mu_{\alpha\beta}(\CandidateMetric)\dot x^\alpha\dot x^\beta=0.
    \label{eq:geodesic}
    \end{equation}
    \item \textbf{Redshift.} For a static chart of the form
    \begin{equation}
    ds_{\mathrm{cand}}^2
    =
    -\Ncand^2(\mathbf{x})\,c^2dt^2
    +\Gcand_{ij}(\mathbf{x})\,dx^i dx^j,
    \label{eq:staticcand}
    \end{equation}
    stationary clocks satisfy
    \begin{equation}
    d\tau=\Ncand(\mathbf{x})\,dt,
    \label{eq:clock}
    \end{equation}
    and stationary light exchange obeys
    \begin{equation}
    \frac{\nu_B}{\nu_A}
    =
    \frac{\Ncand(\mathbf{x}_A)}{\Ncand(\mathbf{x}_B)}.
    \label{eq:redshift}
    \end{equation}
    \item \textbf{Proper time.} Along any timelike worldline,
    \begin{equation}
    d\tau^2
    =
    -\frac{1}{c^2}
    \CandidateMetric_{\mu\nu}dx^\mu dx^\nu.
    \label{eq:propertime}
    \end{equation}
\end{enumerate}
Accordingly, the charge-polarization candidate passes the universal matter/light-coupling gate at current eliminated-candidate scope.
\end{theorem}

\begin{proof}
Propositions \ref{prop:one} and \ref{prop:mattereq} show that all matter sectors remain minimally coupled to \(\CandidateMetric\) after elimination. The Maxwell action \eqref{eq:maxwell} therefore has the standard minimally coupled principal symbol and geometric-optics limit, yielding \eqref{eq:nullcone}. The point-particle action \eqref{eq:pp} varies to the standard geodesic equation \eqref{eq:geodesic}. In a static metric of the form \eqref{eq:staticcand}, stationary observers have \(dx^i=0\), so the line element gives \eqref{eq:clock}; the conserved photon energy along null geodesics then yields \eqref{eq:redshift}. Equation \eqref{eq:propertime} is the timelike line-element definition itself. Since no second metric and no preferred-frame matter term survives by Corollary \ref{cor:nopf}, this is exactly the required Phase E verdict.
\end{proof}

\begin{remark}
Theorem \ref{thm:phaseE} is intentionally phrased after elimination rather than only at architectural action level. The point of the present gate is not merely that the candidate was written with one metric in mind. The point is that, after the response sector is actually solved out, the remaining matter and light laws still reduce to the same one-metric minimal-coupling structure.
\end{remark}

\section{Status Relative to the Remaining Gates}

\begin{proposition}[What this note settles and what remains open]
\label{prop:status}
At the present disciplined scope, the charge-polarization candidate now satisfies the following benchmark-gate tasks.
\begin{enumerate}
    \item The full coupled reduced system has been linearized about an admissible homogeneous benchmark branch.
    \item The observer-accessible infrared spectrum has been derived explicitly, and only the ordinary GR two-tensor sector propagates unsuppressed.
    \item No unsuppressed preferred-frame, scalar, vector, ghost, tachyonic, or second-cone remainder survives at linear infrared order.
    \item The matter and light sector has been re-derived after elimination, and all observer-side null-cone, geodesic, redshift, and proper-time structures continue to run through the single operative metric \(\CandidateMetric\).
\end{enumerate}
What remains open is exactly the next honest work:
\begin{enumerate}
    \item the controlled GR-limit theorem with explicit scaling and decoupling hypotheses;
    \item the residual-sector verdict away from the decisive exterior cancellation regime;
    \item the classification of any remaining higher-order sector as absent, gauge exact, constraint exact, or decoupled.
\end{enumerate}
\end{proposition}

\begin{proof}
Items (1)--(3) are Theorem \ref{thm:phaseD}. Item (4) is Theorem \ref{thm:phaseE}. The remaining items are not proved in the present note and are therefore exactly the honest next burdens left for Phases F and G.
\end{proof}

\section{Conclusion}

The immediate next derivational job after the charge-polarization candidate note was not another packaging pass and not another return to the anchored positive-pair side line. It was the benchmark-gate follow-up for that concrete candidate itself. The present manuscript performs that follow-up at the next honest scope: the full coupled infrared-spectrum test and the universal matter/light-coupling check.

The positive result is precise. On an admissible homogeneous benchmark branch, the full coupled reduced linearization factorizes into the already-admissible GR-compatible base block and an explicit charge-polarization response block that is purely elliptic or algebraic. The observer-accessible infrared poles are therefore exactly the ordinary two tensor poles \(\omega^2=|\mathbf{k}|^2\), and no unsuppressed preferred-frame, scalar, vector, ghost, tachyonic, or second-cone remainder survives at linear infrared order. After elimination, matter and light still couple only through the single operative metric \(\CandidateMetric\), so the null cone, geodesic law, redshift, and proper-time structure remain those of that same one metric.

The scope of that result is still disciplined. This note does not yet prove the controlled GR-limit theorem, and it does not yet classify every residual higher-order sector away from the decisive exterior cancellation regime. Those are now the honest next burdens. But the benchmark-gate sequence is materially advanced: the charge-polarization candidate has now cleared Phase D and Phase E at current branch scope, so the next primary manuscripts should be the controlled-limit note and then the residual-sector verdict, not a return to side-line continuation unless that side line changes the observer equation itself.

\end{document}
